% 星群（综述）
% license CCBYSA3
% type Wiki

本文根据 CC-BY-SA 协议转载翻译自维基百科\href{https://en.wikipedia.org/wiki/Asterism_(astronomy)}{相关文章}。

\begin{figure}[ht]
\centering
\includegraphics[width=8cm]{./figures/c47b677da81e5ec4.png}
\caption{一幅恒星图像，其中有一组显眼的蓝白色亮星。这些亮星共同组成一个名为“衣架”（Coathanger）的星群，形状类似衣架，位于狐狸座（Vulpecula）中。} \label{fig_Asteri_1}
\end{figure}
星群（Asterism）是指在天空中可见的恒星排列或星群图案。星群可以是任何被识别出的恒星形状或组合，因此它的概念比正式定义的 88 个星座更为广泛。星座是以星群为基础的，但与星群不同，星座是具有官方边界的明确区域，整体上覆盖了整个天空。

星群的形态从仅由几颗恒星构成的简单图案，到由大量恒星组成、横跨天空大片区域的复杂结构不等。构成星群的恒星可能是肉眼可见的明亮星，也可能是较暗、甚至需要望远镜才能观测到的星体，但它们通常在亮度上彼此相近。那些较大且较亮的星群对熟悉夜空的人而言具有重要的指引作用。

星群中恒星的排列并不一定反映它们之间的物理联系，而往往是由于观测视角造成的视觉效果。例如，“夏季大三角”仅仅是从地球观察到的一个表观组合，成员恒星之间在空间上并无关联；而“猎户座腰带”中的恒星则都属于猎户座 OB1 星协，“北斗七星”中的七颗星有五颗属于大熊座移动星群（Ursa Major Moving Group）。某些具有真实物理关联的星团，如“毕星团”（Hyades）或“昴星团”（Pleiades），既可以被视为独立的星群，也可能同时是更大星群的一部分。
\subsection{星群与星座的背景}
\begin{figure}[ht]
\centering
\includegraphics[width=8cm]{./figures/590b7d57cdce8a34.png}
\caption{NGC 1664 的图像，其星群呈风筝及尾状，被称为“坎布尔之风筝”。} \label{fig_Asteri_2}
\end{figure}
在许多早期文明中，人们常常将恒星以“连点成线”的方式连接起来，形成类似火柴人形状的图案。其中最早的记录之一来自古印度的《吠檀伽·乔提沙》（Vedanga Jyotisha）以及古巴比伦的天文学文献。\(^\text{[citation needed]}\) 不同的文化辨识出不同的星座，但某些显而易见的恒星排列形式，如猎户座（Orion）和天蝎座（Scorpius），往往会在多个文化的星座体系中同时出现。由于任何人都可以随意组合和命名一组恒星，因此在早期并不存在星座（constellation）与星群（asterism）的明确区别。例如，老普林尼（Pliny the Elder）在其著作《自然史》（Naturalis Historia）中提到了七十二个星群。\(^\text{[3]}\)

包含四十八个星座的一般列表很可能始于天文学家喜帕恰斯（Hipparchus，约公元前190年至公元前120年）。由于当时星座仅被认为是由构成该图形的恒星组成，因此总可以利用剩余的恒星，在既有星座之间创造并插入新的星群。\(^\text{[citation needed]}\)

欧洲人向世界其他地区的探索，使他们接触到了此前从未观测到的恒星。两位以显著扩充南天星座数量而闻名的天文学家是约翰·拜耳（Johann Bayer，1572–1625）与尼古拉·路易·德·拉卡伊（Nicolas Louis de Lacaille，1713–1762）。拜耳列出了十二个图形，这些图形由托勒密（Ptolemy）所无法观测到的南天恒星构成；拉卡伊则在接近南天极的区域创建了十四个新的星群。这些提议中的许多后来被正式接受为星座，而其余部分则保留为星群。\(^\text{[citation needed]}\)

1928年，国际天文学联合会（International Astronomical Union, IAU）根据几何边界将天空精确划分为88个官方星座，将其中包含的所有恒星纳入其范围之内。此后任何新选定的恒星组合或旧有的星座体系，通常都被视为星群。然而，关于“星座”（constellation）与“星群”（asterism）这两个术语之间的技术性区别，至今仍然存在一定的模糊性。\(^\text{[citation needed]}\)
